\chapter{Project Overview}

\section*{i. Document 1}

\subsection*{Purpose of this document}

This System Requirements Document (SRD) defines what the first version of the robot must achieve. It establishes the mission, operating environment, physical envelope, performance targets, functional content, safety baseline, constraints, and measurable acceptance criteria for the V1 prototype. It is the single controlling reference used across mechanical and electrical design, CAD, embedded and navigation software, component selection, procurement, and verification testing. Where later trade-offs become necessary, decisions are recorded against the requirement identifiers defined here.

\subsection*{Project background and rationale}

The project develops a general-purpose mobile robotics platform, beginning with warehouses as the primary target environment. Warehouses are an ideal starting point: they present well-defined logistics challenges, a high proportion of repetitive tasks, and strong, sustained demand for automation. A significant limitation of many existing automation solutions is that they require substantial modification of the facility, installed guidance infrastructure, re-racking, or even a complete redesign of the layout. These changes are costly, disruptive, and slow to deploy.

\subsection*{Solution concept and differentiator}

The core differentiator of this project is a robot that integrates into an existing warehouse without requiring the facility to be redesigned. The base robot is responsible for autonomous navigation, movement, obstacle avoidance, and task execution within the warehouse. The robot is a flexible, modular platform: a common mobile base onto which different modules and attachments can be mounted depending on mission requirements. Version 1 establishes and proves that base platform and its modular interface, on top of which task-specific capability will be added in later versions.

\section*{ii. Document 2}

The project aims to design and develop a modular autonomous mobile robot intended for warehouse environments. The system is conceived as a flexible robotic platform capable of performing basic logistics tasks such as autonomous navigation, obstacle avoidance, and payload transportation.

Unlike traditional single-purpose warehouse machines, the proposed solution focuses on modularity, allowing future integration of additional functional units such as robotic arms, inspection sensors, and specialized handling modules.

Version 1 of the system will focus on the development of a functional mobile base equipped with essential navigation sensors, a payload platform, and basic control capabilities. This first prototype will serve as the foundation for future system expansions and advanced warehouse automation capabilities.
